\chapter{R code}
\label{ch:appendix}

This appendix reports condensed excerpts of the R code that implements the pipeline described in Chapter~\ref{ch:ch3}. The complete pipeline is organised into twenty-one numbered scripts (\texttt{v9/00\_*.R} through \texttt{v9/21\_*.R}), two shared helper modules (\texttt{v9/var\_labels.R}, \texttt{v9/fig\_paths.R}) and one numbers-assembly script (\texttt{v9/20\_thesis\_tables.R}); the full listings are available in the online appendix. The snippets below are self-contained: they rely only on CRAN packages (\texttt{psych}, \texttt{cluster}, \texttt{poLCA}, \texttt{haven}, \texttt{tidyverse}) and require no external toolkit.

\section{A.1\; Rank transformation and z-standardisation}
\label{sec:appA-preprocessing}

Excerpt from \texttt{v9/02\_listwise\_outliers.R} (Italy; the Swedish equivalent lives in \texttt{v9/04\_listwise\_outliers\_sweden.R}):

\begin{lstlisting}[caption={Preprocessing: rank transformation of the 20 non-binary variables and z-standardisation of all 31.}]
# Rank-transform the 20 non-binary variables (average-rank tie-breaking)
italy_ranked <- italy_complete
for (v in nonbinary_vars) {
    italy_ranked[[v]] <- rank(italy_ranked[[v]], ties.method = "average")
}

# Z-standardise all 31 variables (binary + rank-transformed non-binary)
italy_std <- italy_ranked
for (v in all_31) {
    italy_std[[v]] <- as.numeric(scale(italy_std[[v]]))
}
\end{lstlisting}

\section{A.2\; Mahalanobis outlier screen}
\label{sec:appA-outliers}

Excerpt from \texttt{v9/02\_listwise\_outliers.R}:

\begin{lstlisting}[caption={Multivariate outlier detection on the standardised matrix.}]
# Mahalanobis distance on the 31 standardised variables
X       <- as.matrix(italy_std[, all_31])
center  <- colMeans(X)
cov_mat <- cov(X)
mah_dist <- mahalanobis(X, center = center, cov = cov_mat)

# Conservative cut-off: chi-squared with df = 31, p = 0.001
cutoff     <- qchisq(0.999, df = length(all_31))
is_outlier <- mah_dist > cutoff
italy_std  <- italy_std[!is_outlier, ]  # effective removal in v9
\end{lstlisting}

\section{A.3\; Dimension-wise PCA and 6-factor PA-varimax FA}
\label{sec:appA-pca-fa}

Excerpt from \texttt{v9/05\_pca\_fa\_italy.R}:

\begin{lstlisting}[caption={PCA on each of the five conceptual dimensions, then joint FA on the full variable set.}]
# Dimension-wise PCA via eigendecomposition of the correlation matrix
run_pca_block <- function(data_std, vars, block_name, force_n = NULL) {
    X   <- as.matrix(data_std[, vars])
    R   <- cor(X)
    eig <- eigen(R)
    eigenvalues  <- eig$values
    eigenvectors <- eig$vectors
    n_kaiser     <- sum(eigenvalues > 1)
    n_report     <- if (!is.null(force_n)) force_n else n_kaiser
    list(
        block       = block_name,
        eigenvalues = eigenvalues,
        n_retained  = n_report
    )
}

pca_health   <- run_pca_block(italy_std, health_vars,  "Health")
pca_econ     <- run_pca_block(italy_std, econ_vars,    "Economic")
pca_digital  <- run_pca_block(italy_std, digital_vars, "Digital/Social", force_n = 4)
pca_cog      <- run_pca_block(italy_std, cog_vars,     "Cognitive")
pca_subj     <- run_pca_block(italy_std, subj_vars,    "Subjective")

# 6-factor PA-varimax FA on the 29 Italian active variables
library(psych)
X_fa      <- as.matrix(italy_std[, italy_active])
fa_result <- fa(X_fa, nfactors = 6, rotate = "varimax", fm = "pa",
                scores = "regression")
print(fa_result$loadings, cutoff = 0.30, sort = TRUE)
\end{lstlisting}

\section{A.4\; K-means clustering on the standardised variables}
\label{sec:appA-kmeans}

Excerpt from \texttt{v9/07\_cluster\_italy.R}:

\begin{lstlisting}[caption={K-means with one hundred random restarts on the 29 standardised Italian active variables.}]
# Reproducible K-means with 100 random restarts and 500 iterations
X <- as.matrix(italy_std[, active_vars])
set.seed(42)
km6 <- kmeans(X, centers = 6, nstart = 100, iter.max = 500)

# Italy: no micro-cluster smaller than 5% emerges in v9 -- refit with k = 5
set.seed(42)
km5 <- kmeans(X, centers = 5, nstart = 100, iter.max = 500)

# Attach profile labels (order determined by cluster means on CASP)
italy_clust$profile <- factor(
    italy_clust$cluster,
    levels = order_by_casp,
    labels = c("Fragile Resigned", "Fragile Depressed", "Moderate Isolated",
               "Traditional Social", "Connected Active")
)
\end{lstlisting}

\section{A.5\; Latent Class Analysis}
\label{sec:appA-lca}

Excerpt from \texttt{v9/10\_lca.R}:

\begin{lstlisting}[caption={LCA with poLCA; tertile discretisation of non-binary variables; fit sequence for $k = 2, \ldots, 8$.}]
library(poLCA)

fit_lca_sequence <- function(df, vars, k_range, country_label,
                             max_iter = 1000, nrep = 5) {
    formula_lca <- make_formula(vars)      # ~ cbind(v1, v2, ...) ~ 1
    df_complete <- df[complete.cases(df[, vars]), ]

    for (k in k_range) {
        set.seed(42)
        fit <- poLCA(formula_lca, data = df_complete, nclass = k,
                     maxiter = max_iter, nrep = nrep,
                     verbose = FALSE, na.rm = TRUE)
        # Normalised entropy: 1 - E / (N log K), where E = -sum p log p
        entropy <- normalised_entropy(fit$posterior)
        bic_k   <- fit$bic
    }
}
\end{lstlisting}

\section{A.6\; Numbers-for-thesis assembler}
\label{sec:appA-numbers}

Excerpt from \texttt{v9/20\_thesis\_tables.R}:

\begin{lstlisting}[caption={Profile-level summary exported to \texttt{thesis\_numbers.md}.}]
# Profile sizes and cluster means on key variables
key_vars <- c("age", "yedu", "internet", "casp", "loneliness",
              "hope_future", "sn_size_w9", "lifesat")

tab_it <- italy_clust %>%
    filter(!is.na(profile)) %>%
    group_by(profile) %>%
    summarise(
        n   = n(),
        pct = 100 * n() / nrow(italy_clust),
        across(all_of(key_vars), ~ mean(as.numeric(zap_labels(.x)), na.rm = TRUE)),
        .groups = "drop"
    )

writeLines(knitr::kable(tab_it, format = "markdown"),
           "v9/outputs/thesis_numbers.md")
\end{lstlisting}
